% !PRINTED 593
\weberpage{593}
\begin{center}
\S.\,162.\\[4pt]
\gesperrt{The Partial Fundamental Ideal.}
\end{center}

Now let
\weq{1}{\tau = t_1\omega_1 + t_2\omega_2 + \cdots + t_{mn}\omega_{mn}}
be a basis form of $\mathfrak{o}$. Then there will be certain substitutions $\chi$ in $\Phi$, and among them always the identical one, by which the congruence
\weq{2}{\tau \,|\, \chi \equiv \tau \pmod{\mathfrak{p}}}
is satisfied. All these substitutions form a group $X$ contained in $\Phi$, which is also a divisor of the group $\Psi$ to which $\mathfrak{p}$ belongs. For from $\tau$ one obtains (by \S.\,146) a basis form of $\mathfrak{p}$ if one puts for the variables $t$ certain linear functions of new variables with integral rational coefficients; if then $\pi$ is such a basis form of $\mathfrak{p}$, it follows from (2) that $\pi \,|\, \chi$ is divisible by $\pi$, and hence also that $\mathfrak{p} \,|\, \chi = \mathfrak{p}$.

Now $\tau$ satisfies an equation of degree $n$, $f(t) = 0$, whose coefficients are integral functions in $R$ with the variables $t_1, t_2 \ldots$; and if $t$ denotes a new variable, and $\tau_1, \tau_2, \ldots, \tau_n$ are the forms conjugate to $\tau$ with respect to $R$, we have
\weq{3}{f(t, t_1, t_2, \ldots) = f(t) = (t-\tau_1)(t-\tau_2) \ldots (t-\tau_n).}
From this we now obtain, if we repeatedly raise to the power $p$ and apply to the numerical coefficients of $f$ the formula \S.\,161, (1):
\weq{4}{f(t^P, t_1^P, t_2^P, \ldots) \equiv [f(t)]^P \equiv (t^P-\tau_1')(t^P-\tau_2') \ldots (t^P-\tau_n') \pmod{\mathfrak{p}},}
where $\tau_1', \tau_2', \ldots, \tau_n'$ are derived from $\tau_1, \tau_2, \ldots, \tau_n$ by replacing the variables $t_1, t_2, \ldots$ by $t_1^P, t_2^P, \ldots$.

If we now put $t = \tau$ in (4), then $f(t)$ vanishes, and it follows that one of the factors of the last product must be divisible by $\mathfrak{p}$. But this can be expressed by saying that there is in $\Phi$ a substitution $\psi_0$ by which the condition
\weq{5}{\tau^P \equiv \tau' \,|\, \psi_0 \pmod{\mathfrak{p}}}
is satisfied.

From $\tau$ all the integers in $\Omega$ arise, when one puts integral rational numbers for the variables. If one then applies
%CONT>>
